% Page margin: Inside min. 25mm
% Outside margin: Correction margin 50mm,
% Side notes may protrude into the 50mm margin.

According to DIN~5008, the page margins should be min. \qty{2.5}{\centi\metre}.
For work requiring a correction margin on the outside, the correction margin
should be \qty{5}{\centi\metre} wide.

The specifications for page geometry are referred to in print language as the
"type area". In the course of the history of letterpress printing many rules
for the design of a good type area have become established. If you search for
the keywords "type area" and "LaTeX", you will find many papers. However, if in
doubt, stick to the the lecturer's instructions, even if these contradict the
classical rules of rules of type area construction.

The setting of the sentence mirror is done with the package~\ctanref{geometry}.
Listing~\ref{lst:tpl:margin} contains an example of how to use the package.
The option \texttt{includeheadfoot} says that the top or bottom margin should extend to
the top of the header or to the bottom of the footer. Otherwise, the margin is
measured to the body of the text.

\lstinputlisting[language={[LaTeX]{TeX}},style=block,float,caption={set page margins},label={lst:tpl:margin}]{code/tpl_margin.en.tex}

The package~\texttt{geometry} offers many more possibilities that are beyond the scope
of this introduction. which would go beyond the scope of this introduction.
Please refer to the package documentation.


\begingroup
\locbox\evenlayout
\locbox\oddlayout

\printheadingsfalse
\printparametersfalse
%\drawdimensionstrue
\currentpage
\setlayoutscale{0.45}

\oddpagelayoutfalse\savebox{\evenlayout}{{\hsize=0.46\paperwidth\vbox{\centering\footnotesize\pagedesign}}}
\oddpagelayouttrue\savebox{\oddlayout}{{\hsize=0.46\paperwidth\vbox{\centering\footnotesize\pagedesign}}}

\begin{sidewaysfigure*}[p]
\begin{minipage}{0.5\linewidth}
\centering
\usebox{\evenlayout}
\end{minipage}%
\begin{minipage}{0.5\linewidth}
\centering
\usebox{\oddlayout}
\end{minipage}
\caption{Setmirror of this document}
\label{fig:tpl:margin_layout}
\end{sidewaysfigure*}
\endgroup

By inserting the package~\ctanref{showframe} you can display the type area of
your document in a form comparable to Figure~\ref{fig:tpl:margin_layout}.

%\lstinputlisting[language={[LaTeX]{TeX}},style=block,float,caption={displaying the type area},label={lst:tpl:margin_showframe}]{code/tpl_margin_showframe.en.tex}
